\documentclass[11pt]{article}
\usepackage[T1]{fontenc}
\usepackage[utf8]{inputenc}
\usepackage{lmodern,amsmath,amssymb,amsthm,mathtools}
\usepackage[margin=25mm]{geometry}
\usepackage[hidelinks]{hyperref}
\newcommand{\C}{\mathbb C}
\newcommand{\R}{\mathbb R}
\newcommand{\PP}{\mathcal P}
\newcommand{\Fcal}{\mathscr F}
\newcommand{\Bcal}{\mathcal B}
\newcommand{\norm}[1]{\left\lVert#1\right\rVert}
\DeclareMathOperator{\Tr}{Tr}
\DeclareMathOperator{\im}{im}
\DeclareMathOperator{\Gram}{Gram}
\DeclareMathOperator{\cond}{cond}
\title{Growing marked flags in the original arithmetic quotient metrics}
\author{Split-Zero mathematical continuation}
\date{19 September 2026}
\begin{document}
\maketitle

\begin{abstract}
For the original arithmetic source and its specified centre-filled
comparison, a linearly growing marked flag has large positive
individual logarithmic volume changes at every original cutoff.
Its balanced four-cutoff return nevertheless has an explicit
exponentially small bound.  The proof estimates every admissible
relation and terminal representative before taking the original
minima; it retains the full terminal--polynomial cross block.
The same estimates give linearly many large generalized
eigenvalues and a complementary subspace forcing all but
\(O_h(q/\log(q/k))\) eigenvalues into an exponentially short
interval about one.  Exact quotient determinant factorization
then carries the complete signed central response into the
attained complement of the marked flag.  The stated complement
is retained as its actual metric quotient.
\end{abstract}

\section{Original source, complete roots and fixed minimum maps}

Keep the complete hypothetical actual quartet and its full order:
\begin{equation}\tag{GMF1}
\begin{gathered}
 \rho=\tfrac12+\delta+i\gamma,\quad
 0<\delta<\tfrac12,\quad \gamma>2,\quad m_0\ge1,\\
 h(s)=\prod_{\zeta\in\{\rho,\bar\rho,1-\rho,1-\bar\rho\}}
                       (s-\zeta)^{m_0},\quad
 v_h=2\xi/h,\quad
 w_h(y)=|v_h(1/2+iy)|^2/(2\pi),\quad m_k=w_h^{*k},\\
 k\equiv1\pmod4,\ k\ge5,\quad
 e=1+k(m_0-1),\quad q=e(k+1)^2,\quad c=k/2,\\
 \chi(S)=\prod_{a,b=0}^k
 [S-c-(2a-k)\delta-i(2b-k)\gamma]^e,\qquad
 E=\C[S]/(\chi).
\end{gathered}
\end{equation}
The arithmetic source is the unchanged \(k\)-fold convolution,
with no numerical off-critical zero substituted for this original
counterfactual.  All limits fix this complete packet.

Write
\begin{equation}\tag{GMF2}
\begin{gathered}
 \alpha=\pi/2,\qquad
 M(\theta)=\int_\R e^{\theta y}w_h(y)\,dy,\quad
 M_0=M(0),\quad M_\alpha=M(\alpha),\\
 b=(\log M)'(\alpha)>0,\quad T=kb,\qquad
 \sigma(y)=|\Gamma(1/4+iy/2)|^2/(2\pi),\quad
 M_\sigma=\int_\R\sigma(y)\,dy=\sqrt{2\pi},\\
 \widetilde m_k(y)=
 \begin{cases}M_\alpha^k\sigma(y),&|y|\le T,\\
 m_k(y),&|y|>T.\end{cases}
\end{gathered}
\end{equation}
Both complete source masses are retained in every comparison.
This is precisely the central comparison in CPT4:
\(m_k(y)dy=M_\alpha^k\,d\nu_{1,1}\) and
\(\widetilde m_k(y)dy=M_\alpha^k\,d\nu_{0,1}\)
\cite{cpt}.  These equalities identify the two source
presentations without changing their norms.

For a positive source density \(a(y)\), let \(H_N[a]\) be the
Gram of \(1,S,\ldots,S^N\) in the original coordinate \(S=c+iy\).
Its inner product is conjugate-linear in the first variable.
For \(N\ge q-1\), let
\begin{equation}\tag{GMF3}
\begin{gathered}
 J_N:\C[S]_{\le N}\twoheadrightarrow E,\quad
                  J_NP=[P]_\chi,\qquad
 \ker J_N=\chi\C[S]_{\le N-q},\\
 G_N[a]=(J_NH_N[a]^{-1}J_N^*)^{-1},\quad
 R_N[a]=H_N[a]^{-1}J_N^*G_N[a].
\end{gathered}
\end{equation}
A negative degree means the zero space.  Multiplication gives
\(J_NR_N=I_E\), and \(R_Nv\) is orthogonal to every vector in
\(\ker J_N\).  Every representative is \(R_Nv+z\), so its squared
norm is \(v^*G_Nv+\norm z_a^2\).
This proves attainment of the original minimum and the exact
formula (GMF3).

For all claims about cutoff return take
\[
 N\in\{q-1,q,2q-1,2q\},
\]
and the proofs below in fact hold for every integer
\(q-1\le N\le2q\).
For a function \(F_N\) of these cutoffs, retain the signs
\[
 \mathcal R_{\rm ar}F_N=F_{q-1}+F_q-F_{2q-1}-F_{2q},\qquad
 \Bcal_k[a]=\mathcal R_{\rm ar}\log\det G_N[a].
\]
No cutoff is identified with another.

Fix an original root \(\lambda\) and set
\begin{equation}\tag{GMF4}
 \chi_\lambda=\frac{\chi}{(S-\lambda)^e},\qquad
 b_\lambda=\chi_\lambda(\lambda)^{-1},\qquad
 v_\lambda=\left[b_\lambda\frac{\chi}{S-\lambda}\right]_\chi .
\end{equation}
The complete-primary coefficient is nonzero.
It equals \(1/\chi'(\lambda)\) only when \(e=1\).
For \(0\le d\le q-2\), define the fixed flag and its displayed frame
\begin{equation}\tag{GMF5}
 \Fcal_{\lambda,d}
 =\operatorname{span}(v_\lambda,[1]_\chi,\ldots,[S^d]_\chi),
 \qquad F_{\lambda,d}=(v_\lambda,1,S,\ldots,S^d).
\end{equation}
The dimension is \(d+2\): the unique representative of
\(v_\lambda\) has degree \(q-1>d\) with leading coefficient
\(b_\lambda\ne0\), and a degree-\(<q\) representative is zero
in \(E\) only if it is the zero polynomial.

\section{Exact low-polynomial evaluation with the original Gamma mass}

Use the fixed auxiliary constants
\begin{equation}\tag{GMF6}
 \theta_*=\pi/4,\qquad
 r_*=\tanh(\pi/8),\qquad
 \ell_*=-\log r_*=\log\coth(\pi/8),\qquad
 \omega_*=\alpha-\theta_*=\pi/4.
\end{equation}
They are bounds on the original source, with no new source order.
The monic Gamma polynomials and their full norms are
\[
 \sum_{n\ge0}\frac{p_n(y)}{n!}z^n
 =(1+z^2)^{-1/4}e^{y\arctan z},\qquad
 \gamma_n=M_\sigma n!(1/2)_n.
\]
The generating function is the original
Meixner--Pollaczek specialization
\(p_n(y)=n!P_n^{(1/4)}(y/2;\pi/2)\).
Substitution in the generating formula
\cite[(18.23.7)]{dlmfMPgen} gives the displayed expression:
\[
 (1-iz)^{-1/4+iy/2}(1+iz)^{-1/4-iy/2}
       =(1+z^2)^{-1/4}e^{y\arctan z}.
\]
The analytic logarithms are the branches equal to zero at \(z=0\).
The Gamma norm is the original GTM7/CPT8 norm
\cite{providers,cpt}, since
\((1/2)_n=(2n)!/(4^n n!)\).

For \(|z|=r_*<1\), the power series for \(\arctan z\) gives
\(|\arctan z|\le\operatorname{artanh}r_*=\pi/8\), and
\(|1+z^2|\ge1-r_*^2\).
Cauchy's coefficient formula therefore gives
\[
 \frac{|p_n(y)|}{n!}
 \le r_*^{-n}(1-r_*^2)^{-1/4}
                      e^{|y|\operatorname{artanh}r_*}.
\]
The elementary central-binomial estimate
\[
 \frac{n!}{(1/2)_n}=\frac{4^n}{\binom{2n}{n}}\le2n+1
\]
follows because the largest of the \(2n+1\) binomial coefficients
is at least their average.
Thus the complete degree-\(d\) evaluation kernel obeys
\begin{equation}\tag{GMF7}
 K_d^\sigma(y,y)=\sum_{n=0}^d\frac{|p_n(y)|^2}{\gamma_n}
 \le A_de^{\theta_*|y|},\qquad
 A_d=\frac{(d+1)^2 e^{2d\ell_*}}
               {M_\sigma\sqrt{1-r_*^2}}.
\end{equation}
The factor \((d+1)^2\) is the exact sum
\(\sum_{n=0}^d(2n+1)\); every coefficient is included.
Expanding any polynomial in the orthogonal basis and applying
Cauchy--Schwarz proves
\[
 |P(c+iy)|^2\le K_d^\sigma(y,y)\norm P_\sigma^2
 \quad(P\in\C[S]_{\le d}).
\]
The coefficient map \(P(S)\mapsto P(c+iy)\) is invertible,
with determinant \(i^{d(d+1)/2}\) of modulus one, so this
estimate and all source determinant identities concern the
original coefficient space.

Evenness and the exact convolution transform give
\[
 \int_\R e^{\theta|y|}m_k(y)dy
 \le \int(e^{\theta y}+e^{-\theta y})m_k(y)dy
 =2M(\theta)^k,\qquad 0\le\theta\le\alpha.
\]
Combine this identity with (GMF7), first on the entire line
and then on the exact complement of \([-Y,Y]\), to obtain
\begin{equation}\tag{GMF8}
\begin{gathered}
 \norm P_{m_k}^2
 \le2A_dM(\theta_*)^k\norm P_\sigma^2,\\
 \int_{|y|>Y}|P(c+iy)|^2m_k(y)dy
 \le2A_dM_\alpha^k e^{-\omega_*Y}\norm P_\sigma^2
 \quad(Y>0).
\end{gathered}
\end{equation}
For the second inequality use
\(e^{\theta_*|y|}\le e^{-\omega_*Y}e^{\alpha|y|}\)
on the retained exterior.

Let \(C_\Gamma=\sqrt{2\sqrt5}e^{2/3}\) be the original upper
Gamma constant in RTM2/CPT3 \cite{providers,cpt}, and set
\begin{equation}\tag{GMF9}
 \eta_{k,d}=
 \frac{2C_\Gamma A_d}{\omega_*\sqrt{1+T}}e^{-\omega_*T},
 \qquad
 \zeta_{k,d}=2A_de^{-\omega_*T}.
\end{equation}
Indeed
\[
 \int_{|y|>T}|P(c+iy)|^2d\sigma
 \le2C_\Gamma A_d\norm P_\sigma^2
           \int_T^\infty\frac{e^{-\omega_*y}}{\sqrt{1+y}}dy
 \le\eta_{k,d}\norm P_\sigma^2.
\]
For the centre-filled norm keep both of its real regions:
\begin{equation}\tag{GMF10}
 M_\alpha^k(1-\eta_{k,d})\norm P_\sigma^2
 \le\norm P_{\widetilde m_k}^2
 \le M_\alpha^k(1+\zeta_{k,d})\norm P_\sigma^2.
\end{equation}
The lower bound retains the full central Gamma integral.
The upper bound retains its full-line upper bound and the
arithmetic exterior estimate (GMF8) at \(Y=T\).

Put
\begin{equation}\tag{GMF11}
 a_h=\log\frac{M_\alpha}{M(\theta_*)}>0,\qquad
 d_k=\left\lfloor\frac{a_h k}{4\ell_*}\right\rfloor.
\end{equation}
Strict positivity and strict convexity of \(\log M\) follow
by differentiating the actual source integral: its second
derivative is the variance of the positive tilted density
divided by its full mass, and that density is not supported
at a point.  Evenness gives \((\log M)'(0)=0\).
Consequently
\[
 0<a_h=\int_{\theta_*}^{\alpha}(\log M)'(\theta)d\theta
       <\omega_*b.
\]
For \(0\le d\le d_k\), \(2d\ell_*\le a_hk/2\) and
\(\log A_d\le a_hk/2+O_h(\log k)\).
Equations (GMF9) therefore prove, uniformly in these degrees,
\begin{equation}\tag{GMF12}
 \eta_{k,d},\ \zeta_{k,d}
 \le\exp[-a_hk/2+O_h(\log k)].
\end{equation}
The finite inequalities (GMF7)--(GMF10) remain the bounds used
at each stated \(k\); eventual notation has not replaced a guard.

\section{Uniform localization of every high-fibre polynomial}

Write \(B_h=42+8m_0\) and
\(\vartheta_h=\int_{-1}^1w_h(y)dy>0\).
The original threefold convolution lower bound is
\[
 w_h^{*3}(y)\ge c_h e^{-\alpha|y|}(1+|y|)^{-B_h}.
\]
In its convolution with the remaining \(k-3\) factors,
restrict every additional variable to \([-1,1]\).
The sum of those variables has modulus at most \(k-3\);
their retained product mass is \(\vartheta_h^{k-3}\).
The triangle inequality in the exponential and the polynomial
factor therefore proves on the entire real line
\begin{equation}\tag{GMF13}
 m_k(y)\ge c_h\vartheta_h^{k-3}
 e^{-\alpha(|y|+k-3)}(k-2+|y|)^{-B_h}.
\end{equation}
This derivation retains the original numerator and its
convolution mass.

Set
\begin{equation}\tag{GMF14}
\begin{gathered}
 R_k=k\sqrt{\delta^2+\gamma^2},\quad Y=\sqrt{kq},\\
 T\le Y,\qquad Y\le q/2,\qquad Y+2R_k<q,\\
 \underline m_{k,q}
 =c_h\vartheta_h^{k-3}e^{-\alpha(2q+k-3)}
                        (k-2+2q)^{-B_h},\\
 \rho_Y=\frac{Y+R_k}{q-R_k}<1,\qquad
 \mathcal M_k=2M_0^k+M_\alpha^kM_\sigma.
\end{gathered}
\end{equation}
The three displayed inequalities are explicit finite guards.
They hold eventually at each fixed packet, since \(q\ge(k+1)^2\).
The number \(\underline m_{k,q}\) is a lower bound for the
actual density at every point of \([q,2q]\).
Also
\(\int(m_k+\widetilde m_k)\le\mathcal M_k\), by (GMF2) and
\(\int m_k=M_0^k\).

Let \(\Psi(y)\) be monic of degree \(q\) or \(q-1\), with
every root, counted with multiplicity, of modulus at most \(R_k\).
For every \(F=\Psi p\) of degree at most \(2q\), define
\begin{equation}\tag{GMF15}
 \beta_k=
 \frac{\mathcal M_k}{\underline m_{k,q}}
 \frac{(q+2)^2}{q}\,9^{2(q+1)}
                     \rho_Y^{2(q-1)}.
\end{equation}
Then
\begin{equation}\tag{GMF16}
 \int_{|y|\le Y}|F(y)|^2(m_k(y)+\widetilde m_k(y))dy
 \le\beta_k\norm F_{m_k}^2,\qquad
 (1-\beta_k)\norm F_{m_k}^2
 \le\norm F_{\widetilde m_k}^2
 \le(1+\beta_k)\norm F_{m_k}^2.
\end{equation}

To prove the full coefficient estimate, \(\deg p\le q+1\).
The shifted orthonormal Legendre polynomials on \([q,2q]\)
are
\[
 q^{-1/2}\sqrt{2j+1}\,P_j(2y/q-3),\qquad 0\le j\le q+1.
\]
Use the orthogonality and recurrence in
\cite[Tables 18.3.1 and 18.9.1]{dlmfLegendre}.
For \(|x|\le4\), their recurrence, \(P_0=1,P_1=x\), gives
\(|P_j(x)|\le9^j\), because
\[
 8\cdot9^j+9^{j-1}<9^{j+1}.
\]
Under \(Y\le q/2\), \(|2y/q-3|\le4\) on \(|y|\le Y\).
Expand the complete polynomial \(p\) in this basis and
apply Cauchy--Schwarz.  Summing every evaluation square yields
\[
 \sup_{|y|\le Y}|p(y)|^2
 \le\frac{(q+2)^2}{q}9^{2(q+1)}
                         \int_q^{2q}|p(x)|^2dx.
\]
All roots retained in \(\Psi\) give
\[
 |\Psi(y)|\le(Y+R_k)^{\deg\Psi}\quad(|y|\le Y),\qquad
 |\Psi(x)|\ge(q-R_k)^{\deg\Psi}\quad(q\le x\le2q).
\]
Use the full mass \(\mathcal M_k\) for the numerator and the
actual lower density \(\underline m_{k,q}\) for the indicated
part of the denominator.  Division for \(p\ne0\) gives a
factor \(\rho_Y^{2\deg\Psi}\), at most
\(\rho_Y^{2(q-1)}\).  The zero polynomial is immediate.
This proves the first inequality (GMF16).
The measures agree for \(|y|>T\), and \(T\le Y\); hence
\[
 |\norm F_{\widetilde m_k}^2-\norm F_{m_k}^2|
 \le\int_{|y|\le Y}|F|^2(m_k+\widetilde m_k).
\]
This proves its two relative inequalities without selecting
a preferred representative.

In the original coefficient coordinate, the complete relation
polynomials have the form \(\chi P\), and every representative
of \(v_\lambda\) has the form
\begin{equation}\tag{GMF17}
 b_\lambda\frac{\chi}{S-\lambda}+\chi P
 =\frac{\chi}{S-\lambda}
                  [b_\lambda+(S-\lambda)P].
\end{equation}
The real-coordinate monic factors are
\(i^{-q}\chi(c+iy)\) and
\(i^{-(q-1)}(\chi/(S-\lambda))(c+iy)\).
Their roots have modulus at most \(R_k\).
The factors \(i^q,i^{q-1}\), and \(b_\lambda\) in (GMF17)
are retained scalars; (GMF16) applies to all coefficients and
therefore applies before both original minima.
Writing
\(E_{N,\lambda}[a]=v_\lambda^*G_N[a]v_\lambda\), it gives
\begin{equation}\tag{GMF18}
 1-\beta_k\le
 \frac{E_{N,\lambda}[\widetilde m_k]}{E_{N,\lambda}[m_k]}
 \le1+\beta_k,\qquad q-1\le N\le2q.
\end{equation}
At \(N=q-1\) the fibre consists of the unique degree-\(<q\)
representative; the same inequalities still hold.

Taking logarithms of the explicitly positive bound (GMF15),
and using
\[
 \log\rho_Y
 =-\tfrac12\log(q/k)
   +\log(1+R_k/\sqrt{kq})-\log(1-R_k/q),
\]
gives the fixed-packet estimate
\begin{equation}\tag{GMF19}
 \log\beta_k
 \le-(q-1)\log(q/k)+O_h(q+k+\log q).
\end{equation}
In particular \(\beta_k\) tends to zero faster than
\(\exp(-c q)\) for every fixed \(c>0\).
This does not remove its finite value from the following guards.

\section{Original attained low metric and the entire cross block}

The original global envelope and polynomial norm comparison are
\begin{equation}\tag{GMF20}
\begin{gathered}
 m_k(y)\ge a_k(1+y^2)^{-\beta_0}\sigma(y),\qquad
 \beta_0=21+4m_0,\\
 a_k=\frac{c_h\vartheta_h^{k-3}e^{-\alpha(k-3)}(k-2)^{-B_h}}
                 {C_\Gamma 2^{B_h/2}},\qquad
 D_d=[1+4(d+\beta_0)^2]^{\beta_0},\\
 \norm P_{m_k}^2\ge a_kD_d^{-1}\norm P_\sigma^2
                    \qquad(\deg P\le d).
\end{gathered}
\end{equation}
These are the original GTM1 and ACW25--26 bounds
\cite{providers,centralProvider}.  They also follow from
the sharper CPT7 bound
\([1+4(d+1)^2]^{-\beta_0}\), since \(\beta_0\ge1\)
\cite{cpt}.  Thus the incoming notation \(a_k^{\rm AMT}\)
in this calculation is fixed as this actual \(a_k\); no
independent source mass is introduced.

Use (GMF8),(GMF10),(GMF20) to define the finite low exterior
fractions
\begin{equation}\tag{GMF21}
\begin{gathered}
 \ell_m=
 2A_dM_\alpha^k\frac{D_d}{a_k}e^{-\omega_*Y},\qquad
 \ell_f=\frac{2A_d}{1-\eta_{k,d}}e^{-\omega_*Y},\\
 e_m=(\sqrt{\beta_k}+\sqrt{\ell_m})^2,\qquad
 e_f=\left(\sqrt{\frac{\beta_k}{1-\beta_k}}
                              +\sqrt{\ell_f}\right)^2,\\
 \eta_{k,d}<1/2,\quad \beta_k<1/4,\quad
 e_m<1/4,\quad e_f<1/4,\qquad
 z_m=\frac{e_m}{1-e_m},\quad z_f=\frac{e_f}{1-e_f}.
\end{gathered}
\end{equation}
The subscripts \(m,f\) denote the actual arithmetic and
centre-filled sources respectively.
Each displayed strict inequality is a finite guard.

For a low polynomial \(P\) of degree at most \(d\), (GMF21)
bounds its fraction of squared norm outside \([-Y,Y]\).
For any high polynomial covered by (GMF16), that formula
bounds its central fraction in the original source by
\(\beta_k\), and in the filled source by
\(\beta_k/(1-\beta_k)\).
The latter uses
\(\norm F_{\widetilde m_k}^2\ge(1-\beta_k)\norm F_{m_k}^2\).
Split the pairing into the two \emph{complete} regions
\(|y|\le Y\) and \(|y|>Y\), apply Cauchy--Schwarz on each,
and add before squaring.  This proves
\begin{equation}\tag{GMF22}
 |\langle P,F\rangle_a|^2
 \le e_a\norm P_a^2\norm F_a^2,\qquad
 a=m_k,\widetilde m_k,
\end{equation}
where \(e_{m_k}=e_m,e_{\widetilde m_k}=e_f\).
This bound holds for every high polynomial, uniformly in all
its coefficients.

Let \(\Pi_{\mathcal L_N}^a\) be the orthogonal projection in
\(\C[S]_{\le N}\) onto the \emph{whole} relation space
\(\mathcal L_N=\chi\C[S]_{\le N-q}\).
The original attained representative of a low class is
\(P-\Pi_{\mathcal L_N}^aP\).
Every vector of \(\mathcal L_N\) is a high polynomial in
(GMF22).  Taking the supremum over its unit sphere proves
\(\norm{\Pi_{\mathcal L_N}^aP}_a^2\le e_a\norm P_a^2\).
Pythagoras then gives the full matrix statement in the
fixed low-monomial frame:
\begin{equation}\tag{GMF23}
 (1-e_a)H_d[a]\preceq G_N[a]\big|_{\PP_d}
                         \preceq H_d[a].
\end{equation}
For the empty relation space at \(N=q-1\), this is the
exact identity \(G_{q-1}|_{\PP_d}=H_d\) with an admissible
looser bound.  Thus the original minimum has been evaluated
up to a proved error for every relation direction, rather
than by its value on one trial representative.

Let \(v_N=R_N[a]v_\lambda\) be the terminal minimizer.
By (GMF17) it is a high polynomial.  Its orthogonality to
\(\mathcal L_N\) gives
\[
 \langle v_N,P-\Pi_{\mathcal L_N}^aP\rangle_a
                       =\langle v_N,P\rangle_a.
\]
Equation (GMF22) and the lower inequality of (GMF23)
therefore bound the squared angle of the terminal value
with the \emph{entire} low quotient subspace by \(z_a\).
Explicitly, in the flag frame the Gram is
\[
 \begin{pmatrix} E_{N,\lambda}[a]&B_{N,a}\\
                 B_{N,a}^*&C_{N,a}\end{pmatrix},\qquad
 C_{N,a}=G_N[a]|_{\PP_d}.
\]
The coefficient row \(B_{N,a}\) contains all terminal--low
pairings.  Completing the square gives
\begin{equation}\tag{GMF24}
\begin{gathered}
 j_{N,a}=
 \frac{B_{N,a}C_{N,a}^{-1}B_{N,a}^*}{E_{N,\lambda}[a]},
 \quad 0\le j_{N,a}\le z_a<1,\\
 \det\Gram_{G_N[a]}F_{\lambda,d}
 =E_{N,\lambda}[a]\det C_{N,a}\,(1-j_{N,a}).
\end{gathered}
\end{equation}
For the upper bound on \(j\), its numerator divided by \(E\)
is exactly the supremum of the squared normalized terminal
pairing over the entire low subspace, by Cauchy--Schwarz
in the \(C_{N,a}\)-metric, with equality in the direction
\(C_{N,a}^{-1}B_{N,a}^*\).
Thus no entry of the original cross block has been discarded.

\section{Growing marked determinant and its four-cutoff cancellation}

For \(0\le x<1\) set \(h_0(x)=-\log(1-x)\), and define
\begin{equation}\tag{GMF25}
 \Omega_{k,d}=
 h_0(\beta_k)
 +(d+1)\{h_0(e_m)+h_0(e_f)\}
 +h_0(z_m)+h_0(z_f).
\end{equation}
The preceding finite guards make every term well defined.
Equations (GMF18),(GMF23),(GMF24) imply, simultaneously at
every original cutoff,
\begin{equation}\tag{GMF26}
\begin{split}
 \left|
 \log\frac{\det\Gram_{G_N[\widetilde m_k]}F_{\lambda,d}}
              {\det\Gram_{G_N[m_k]}F_{\lambda,d}}
 -\log\frac{\det H_d[\widetilde m_k]}{\det H_d[m_k]}
 \right|\le\Omega_{k,d}.
\end{split}
\end{equation}
Indeed the terminal norm-ratio logarithm is at most
\(h_0(\beta_k)\) in absolute value.
For each source the logarithmic determinant difference
between \(C_{N,a}\) and \(H_d[a]\) belongs to
\([-(d+1)h_0(e_a),0]\).
The two cross-block logarithms belong respectively to
\([-h_0(z_a),0]\).
Adding their absolute error bounds proves (GMF26) with
the stated full rank \(d+1\).

The source determinant term in (GMF26) is the same at
each cutoff \(N\).  Its four signed coefficients sum to zero.
Consequently
\begin{equation}\tag{GMF27}
 \boxed{
 \left|\mathcal R_{\rm ar}
 \log\frac{\det\Gram_{G_N[\widetilde m_k]}F_{\lambda,d}}
              {\det\Gram_{G_N[m_k]}F_{\lambda,d}}
 \right|\le4\Omega_{k,d}.}
\end{equation}

Take \(0\le d\le\min(d_k,q-2)\).
The finite constants show
\[
 \log\ell_m,\ \log\ell_f
 \le-\omega_*\sqrt{kq}+O_h(k+\log q).
\]
For \(\ell_m\), this uses the explicit formula for \(a_k\)
in (GMF20), the bound \(\log A_d=O_h(k+\log k)\), and
\(\log D_d=O_h(\log q)\); no unrecorded mass is cancelled.
For \(\ell_f\) use \(\eta_{k,d}<1/2\).
The bound (GMF19) for \(\beta_k\) is smaller than these
bounds on the eventual fixed-packet range.
Since \((\sqrt x+\sqrt y)^2\le2x+2y\), each \(e_a\) obeys
the same leading exponential bound.  The functions
\(h_0(x)\) and \(x/(1-x)\) are at most fixed multiples
of \(x\) on the finite guards in (GMF21).
Multiplication by \(d+1=O_h(k)\) adds only \(O_h(\log k)\)
to the logarithmic error.  Hence
\begin{equation}\tag{GMF28}
 \Omega_{k,d}\le
 \exp[-(\pi/4)\sqrt{kq}+O_h(k+\log q)]
 =O_h(e^{-(\pi/8)\sqrt{kq}}).
\end{equation}
All the strict finite guards hold eventually, uniformly in
this growing range, as proved by these explicit bounds.
At fixed packet \(d_k\le q-2\) eventually, since \(d_k=O_h(k)\)
and \(q\ge(k+1)^2\).

The per-cutoff logarithmic volumes themselves have a very
different evaluated lower bound.  From (GMF8),(GMF10),
\begin{equation}\tag{GMF29}
 \log\frac{\det H_d[\widetilde m_k]}{\det H_d[m_k]}
 \ge(d+1)\bigl[
 ka_h-\log(2A_d)+\log(1-\eta_{k,d})\bigr].
\end{equation}
This is the determinant of a full Loewner comparison on
the unchanged degree-\(d\) coefficient space.
At \(d=d_k=a_hk/(4\ell_*)+O(1)\),
\(\log A_d=2d\ell_*+2\log(d+1)+O(1)\).
Inserting this into (GMF29), and then (GMF26), yields
\begin{equation}\tag{GMF30}
 \boxed{
 \log\frac{\det\Gram_{G_N[\widetilde m_k]}F_{\lambda,d_k}}
              {\det\Gram_{G_N[m_k]}F_{\lambda,d_k}}
 \ge\frac{a_h^2}{8\ell_*}k^2-O_h(k\log k)}
\end{equation}
uniformly over the full original cutoff window.
The coefficient uses the fixed actual mass ratio \(a_h\).
It has not been obtained from a rank allocation of the
full signed determinant.

\section{Relative eigenvalues: both the large part and the majority}

Let
\[
 \mathcal A_N=
 G_N[m_k]^{-1/2}G_N[\widetilde m_k]G_N[m_k]^{-1/2}
\]
be the positive Hermitian representative of the relative metric.
The eigenvalues are independent of the chosen fixed value frame,
being the generalized eigenvalues of the two original forms.
On the low class subspace of dimension \(d+1\),
(GMF23),(GMF8),(GMF10) give the exact Rayleigh bound
\begin{equation}\tag{GMF31}
 \frac{\norm{[P]_\chi}_{G_N[\widetilde m_k]}^2}
      {\norm{[P]_\chi}_{G_N[m_k]}^2}
 \ge
 \frac{(1-e_f)(1-\eta_{k,d})e^{ka_h}}{2A_d}
 \quad(0\ne P\in\PP_d).
\end{equation}
The map from this class subspace into the Euclidean space of
\(\mathcal A_N\) is multiplication by \(G_N[m_k]^{1/2}\);
it is an isomorphism onto a \((d+1)\)-dimensional subspace.
If fewer than \(d+1\) eigenvalues were at least the right
side of (GMF31), this subspace would intersect the span
of the remaining smaller-eigenvalue eigenvectors nontrivially,
contradicting (GMF31).
At \(d=d_k\) this proves
\begin{equation}\tag{GMF32}
 \#\{\lambda_j(\mathcal A_N):
 \log\lambda_j\ge a_hk/2-O_h(\log k)\}\ge d_k+1.
\end{equation}
The argument proves the needed finite-dimensional min--max
consequence directly; it does not choose trial eigenvectors
or a preferred quotient section.

For the complementary estimate, take any literal monic
divisor \(D\mid\chi\) of degree \(j\), including its selected
repetitions.  In real coordinates make the monic factor
\(i^{-j}D(c+iy)\); every root still has modulus at most \(R_k\).
For any polynomial of degree at most \(2q\) divisible by
this factor, repeat the full-coefficient Legendre proof of
(GMF16), this time allowing the remaining degree through
\(2q\).  It gives relative error at most
\begin{equation}\tag{GMF33}
 B_{k,j}=A_k^{\rm high}\rho_Y^{2j},\qquad
 A_k^{\rm high}=
 \frac{\mathcal M_k}{\underline m_{k,q}}
 \frac{(2q+1)^2}{q}\,9^{4q}.
\end{equation}
Specifically the Legendre sum is bounded by
\((2q+1)^2 9^{4q}/q\); the remaining steps use exactly
the numerator mass and denominator density in (GMF14).
Put
\begin{equation}\tag{GMF34}
 j_k=\left\lceil
 \frac{\log A_k^{\rm high}+q}{2\log(1/\rho_Y)}
 \right\rceil ,
 \qquad 1\le j_k<q/2.
\end{equation}
The displayed degree inequalities are finite guards.
They hold eventually: \(\log A_k^{\rm high}=O_h(q+k+\log q)\),
while
\(2\log(1/\rho_Y)=\log(q/k)+o_h(1)\to\infty\).
This also proves
\(j_k=O_h(q/\log(q/k))\).
The ceiling in (GMF34) gives \(B_{k,j_k}\le e^{-q}\) exactly.

Choose the divisor \(D\) as the product of the first \(j_k\)
factors in the original ordered factor list for \(\chi\),
with each repetition retained.  The multiplication map
\begin{equation}\tag{GMF35}
 \C[S]/(\chi/D)\longrightarrow E,\qquad
 [P]_{\chi/D}\longmapsto[DP]_\chi
\end{equation}
is well defined and injective: \(\chi\mid DP\) is
equivalent in the polynomial domain to \(\chi/D\mid P\).
Its image \(\mathcal V_D\) has dimension \(q-j_k\).
Every polynomial representative of a class in \(\mathcal V_D\)
is divisible by \(D\), since its difference from a displayed
representative is divisible by \(\chi\).
The bounds (GMF33) therefore apply to every representative
in its original minimum fibre.  Taking both minima proves
\begin{equation}\tag{GMF36}
 (1-e^{-q})\,G_N[m_k]|_{\mathcal V_D}
 \preceq G_N[\widetilde m_k]|_{\mathcal V_D}
 \preceq(1+e^{-q})\,G_N[m_k]|_{\mathcal V_D}.
\end{equation}
This is the same exact divisibility-to-minimum map as
CPT25--CPT27, here with \(g=\chi/D\) and the new enlarged
localization interval \(Y\).

In the relative Euclidean coordinates used above, the
subspace in (GMF36) has codimension \(j_k\).
If the spectral subspace below \(1-e^{-q}\) had dimension
greater than \(j_k\), it would intersect this subspace,
contradicting its lower Rayleigh bound.
The same dimension argument above \(1+e^{-q}\) proves
\begin{equation}\tag{GMF37}
 \boxed{
 \#\{\lambda_j(\mathcal A_N)\notin
          [1-e^{-q},1+e^{-q}]\}
 \le2j_k=O_h(q/\log(q/k)).}
\end{equation}
Consequently, for every bounded continuous function \(f\),
\[
 \left|\frac1q\sum_{j=1}^q f(\lambda_j)-f(1)\right|
 \le\sup_{|x-1|\le e^{-q}}|f(x)-f(1)|
                    +4\norm f_\infty\,j_k/q\longrightarrow0 .
\]
Both estimates are uniform over all original cutoffs.
The linearly many large eigenvalues in (GMF32) coexist
with this concentration because their proportion tends
to zero.  Their logarithmic volume contribution is
explicitly retained in (GMF30).

\section{Marked derivatives, period data and the attained complement}

The original marked coefficient module is
\[
 \mathscr H=\C[u,u^{-1},t][S]\big/
       (u\partial_S+\chi-t)\C[u,u^{-1},t][S].
\]
Its vectors have unique representatives of degree below
\(q\); this quotient is a module, with the original marked
connection, and is not used as a quotient ring.
Indeed subtracting the differential relation applied to
a leading monomial reduces any polynomial's degree.
A nonzero \(P\) has
\(\deg((u\partial_S+\chi-t)P)=q+\deg P\), proving
uniqueness of a representative of degree below \(q\).
The cochain operator \(\partial_t-S/u\) descends because
\[
 [\partial_t-S/u,u\partial_S+\chi-t]=-1+1=0.
\]
On the chosen representatives it is exactly the original
\(\nabla_t=\partial_t-A(t)/u\).
At \(t=0\), identify its coefficient representatives with
the displayed vector space \(E\), and fix \(u\ne0\).
The exact marked recurrence from the marked-map
calculation \cite{markedProvider} follows directly here.
Let \(R_j\) be the unique representative of
\(S^j\chi/(S-\lambda)\).  The cochain formula gives
\(\nabla_t^jv_\lambda=b_\lambda(-u)^{-j}R_j\).
Applying the original relation to \(S^{j-1}\) gives
\[
 R_j-\lambda R_{j-1}
   =tS^{j-1}-u(j-1)S^{j-2}\qquad(1\le j\le q),
\]
whose right side already has degree below \(q\).
The final term is zero at \(j=1\).
Thus at \(t=0\) one has
\[
 v_0=v_\lambda,\quad v_j=(\nabla_t^j v_\lambda)|_{t=0},
 \qquad v_1=-\lambda v_0/u,
\]
\begin{equation}\tag{GMF38}
 v_j+\frac{\lambda}{u}v_{j-1}
   =b_\lambda(j-1)(-u)^{1-j}S^{j-2},
                    \qquad 2\le j\le q.
\end{equation}
It retains the complete-primary \(b_\lambda\).
For \(j=2\), insert the displayed \(v_1\), giving
\(v_2=\lambda^2v_0/u^2-b_\lambda/u\).
For each successive \(j\), (GMF38) expresses \(v_j\)
as a nonzero scalar multiple of \(S^{j-2}\) plus preceding
columns.  Thus in increasing order
\[
 (v_0,v_2,\ldots,v_{d+2})
       =F_{\lambda,d}\,\mathcal T_{\lambda,d}(u),
\]
where the triangular matrix has diagonal
\[
 1,\ b_\lambda(-u)^{-1},\
 2b_\lambda(-u)^{-2},\ldots,
 (d+1)b_\lambda(-u)^{-(d+1)}.
\]
Its complete modulus-squared determinant is
\begin{equation}\tag{GMF39}
 |\det\mathcal T_{\lambda,d}(u)|^2
 =|b_\lambda|^{2(d+1)}((d+1)!)^2
                        |u|^{-(d+1)(d+2)}.
\end{equation}
This proves the partial marked frame map and the exact
Gram congruence.  The scalar cancels between the two source
metrics at each cutoff because the frame is fixed.
Its magnitude is never set to one.  Equations (GMF26)--(GMF30)
therefore hold for the literal marked derivative columns.
The original period isomorphism has its own fixed matrix
\(\Pi\) at the retained nonzero period and branch.  Its exact
transported metric is
\[
 \widehat G_N[a]=(\Pi^{-1})^*G_N[a]\Pi^{-1},\qquad
 (\Pi F_{\lambda,d})^*\widehat G_N[a](\Pi F_{\lambda,d})
                  =F_{\lambda,d}^*G_N[a]F_{\lambda,d}.
\]
Thus the same estimate holds on the literal period images
when equipped with the transported arithmetic metric.
The Euclidean period Gram is separately
\(F_{\lambda,d}^*\Pi^*\Pi F_{\lambda,d}\);
its exact map is the displayed \(\Pi\), and it is not
substituted for either attained arithmetic form here.

Let \(r=d+2\), and complete \(F_{\lambda,d}\) by any fixed
coefficient columns \(C_{\lambda,d}\) to an invertible frame
\(\mathcal J=(F_{\lambda,d},C_{\lambda,d})\) of \(E\).
Use the quotient frame consisting of the images of
\(C_{\lambda,d}\) in \(E/\Fcal_{\lambda,d}\).
Then the actual quotient map \(\pi_\Fcal\) satisfies
\(\pi_\Fcal\mathcal J=(0,I_{q-r})\).
In this completed frame write the original Gram
\[
 \mathcal J^*G_N[a]\mathcal J=
 \begin{pmatrix}A_{N,a}&B_{N,a}\\B_{N,a}^*&C_{N,a}\end{pmatrix}.
\]
Every value \(z\) in the quotient is represented by
\(\mathcal J(x,z)\).  Completing its squared norm gives
the unique minimum at \(x=-A_{N,a}^{-1}B_{N,a}z\).
Consequently its exact attained metric is
\begin{equation}\tag{GMF40}
\begin{gathered}
 Q_{\Fcal,N}[a]
 =(\pi_\Fcal G_N[a]^{-1}\pi_\Fcal^*)^{-1}
 =C_{N,a}-B_{N,a}^*A_{N,a}^{-1}B_{N,a},\\
 |\det\mathcal J|^2\det G_N[a]
   =\det A_{N,a}\det Q_{\Fcal,N}[a].
\end{gathered}
\end{equation}
All original cross entries remain in this Schur complement.
The fixed Jacobian \(|\det\mathcal J|^2\) cancels between
the two sources.

Define the exact central signed endpoint value
\(\delta_k^{\rm cen}=\Bcal_k[m_k]-\Bcal_k[\widetilde m_k]\).
Taking the actual four signs in (GMF40) and using (GMF27)
proves the complete finite receiving theorem
\begin{equation}\tag{GMF41}
\begin{split}
 \mathcal R_{\rm ar}
 \log\frac{\det Q_{\Fcal,N}[\widetilde m_k]}
              {\det Q_{\Fcal,N}[m_k]}
 &=-\delta_k^{\rm cen}+\mathcal E_{\Fcal,k},\\
 \mathcal E_{\Fcal,k}
 &:=-\mathcal R_{\rm ar}
 \log\frac{\det\Gram_{G_N[\widetilde m_k]}F_{\lambda,d}}
              {\det\Gram_{G_N[m_k]}F_{\lambda,d}},
 \qquad |\mathcal E_{\Fcal,k}|\le4\Omega_{k,d}.
\end{split}
\end{equation}
This theorem evaluates the error of transport to the
\emph{actual attained complement}.  Its endpoint value is
the unchanged central endpoint, with the displayed minus
sign.  The accepted finite-amplitude endpoint evaluation is
inserted below with its full source error; no first
derivative is substituted for this endpoint difference.

For clarity the exact map to an independently specified
original observation \(O:E\to Y\) has its own criterion.
It factors through \(\pi_\Fcal\) precisely when
\(O|_\Fcal=0\): the necessity follows by applying the
factorization to the kernel of \(\pi_\Fcal\), and the
sufficiency follows by defining \(\bar O([v])=Ov\).
Without that vanishing the original observation remains
the pair of coefficient maps
\((OF_{\lambda,d},OC_{\lambda,d})\) in the completed frame,
with their full blocks.  In particular (GMF40) has not
identified its quotient with an original observation
quotient or with an induced arithmetic action.

\section{The preceding constant-class receiver is sharpened}

At \(N=q-1\) the relation space is zero.  Define the original
one-packet variance and two retained tails by
\[
 V_0=\frac{\int y^2w_h(y)dy}{M_0},\quad
 D_0=\log(M_\alpha/M_0),\quad
 s_T=\int_{|y|>T}\sigma(y)dy,\quad
 p_{k,T}=M_0^{-k}\int_{|y|>T}m_k(y)dy.
\]
The auxiliary unit-mass law \(w_h/M_0\) has mean zero by
evenness and variance \(V_0>0\).
Its \(k\)-fold sum has variance \(kV_0\), by expansion of
the centred square and vanishing of every cross mean.
Integrating \(y^2/T^2\ge1\) on \(|y|>T=kb\) gives
\[
 0\le p_{k,T}\le V_0/(kb^2),\qquad
 0\le s_T\le
 \frac{2C_\Gamma}{\alpha\sqrt{1+T}}e^{-\alpha T}.
\]
Keeping both tails and both original masses gives the
exact constant-class ratio
\begin{equation}\tag{GMF42}
 \frac{\norm{1}_{G_{q-1}[\widetilde m_k]}^2}
      {\norm{1}_{G_{q-1}[m_k]}^2}
 =e^{kD_0}(M_\sigma-s_T)+p_{k,T}.
\end{equation}
This verifies B7--B8 of the marked arithmetic integration
source \cite{markedProvider}.
The terminal Rayleigh quotient in (GMF18) is at most
\(1+\beta_k\); the ratio of the largest to the smallest
generalized eigenvalue is therefore at least
\begin{equation}\tag{GMF43}
 \cond(\mathcal A_{q-1})
 \ge\frac{e^{kD_0}(M_\sigma-s_T)+p_{k,T}}{1+\beta_k}
 =\sqrt{2\pi}\,e^{kD_0}(1-o_h(1)).
\end{equation}
If the sharper CPT pair radius is used in the terminal
comparison, the same proof uses its exact upper factor
instead; both are original-data bounds.
At every other original cutoff, (GMF23) with \(d=0\)
gives the uniform extension
\[
 \cond(\mathcal A_N)\ge
 \frac{(1-e_f|_{d=0})
       [e^{kD_0}(M_\sigma-s_T)+p_{k,T}]}
      {1+\beta_k}.
\]
Here the original arithmetic constant-class denominator
is at most \(M_0^k\), and the filled numerator is at least
\((1-e_f|_{d=0})\int\widetilde m_k\); the terminal
Rayleigh upper bound remains (GMF18).

For \(d=0\), \(H_0[m_k]=M_0^k\).
Equation (GMF10) therefore gives
\[
 \left|\log\frac{H_0[\widetilde m_k]}{H_0[m_k]}
            -kD_0-\log M_\sigma\right|
 \le h_0(\eta_{k,0})+\log(1+\zeta_{k,0}).
\]
Return it through the complete cross-block formula
(GMF26) and then the marked congruence (GMF39).
This proves at \emph{every} original cutoff the finite bound
\begin{equation}\tag{GMF44}
\begin{split}
 \left|
 \log\frac{\det\Gram_{G_N[\widetilde m_k]}(v_\lambda,\nabla_t^2v_\lambda)}
              {\det\Gram_{G_N[m_k]}(v_\lambda,\nabla_t^2v_\lambda)}
       -kD_0-\tfrac12\log(2\pi)
 \right|
 \le\Omega_{k,0}+h_0(\eta_{k,0})+\log(1+\zeta_{k,0}).
\end{split}
\end{equation}
Its right side is exponentially small at fixed packet.
This strengthens the prior first-cutoff B11 remainder
\(O_h(k^2/q^2)\), using full-fibre localization instead
of only its parity estimate.  The former statement and
its original proof remain preserved.

\section{The evaluated central value in the original attained complement}

The accepted signed arithmetic source proof
\cite[C32 and SAI1--SAI11]{signedCentral} evaluates the
\emph{finite-amplitude endpoint} appearing in (GMF41):
\[
 I_h=\int_0^\alpha[(\log M)'(\theta)]^2\,d\theta>0,\qquad
 \delta_k^{\rm cen}
  =\frac{\log2}{\pi}I_h k^2+E_k^{\rm cen}.
\]
On its original fixed-packet range the retained error is
\[
 |E_k^{\rm cen}|
 \le \mathfrak C_{h,k}^{\rm cen},\qquad
 \mathfrak C_{h,k}^{\rm cen}
 =A_h\left[
 \frac{k^2}{\log^3(q+2)}+k\log^2(q+2)+1+\frac{k^4}{q^2}
 \right].
\]
Here \(A_h\) is the fixed-packet bound in that accepted
source calculation, with its original exact Robin and
Neumann guards.  Specifically its finite monic-band part
is \(\sum_{r=0}^q c_rE_H(q+r)+k^2I_h/(16\pi q^2)\)
from SAI11; its source-localization and smoothing errors
are added as in C3,C8 before the displayed fixed-packet
bound is taken.  These terms are inherited with their
whole-line integration and both original endpoint weights.
This use is a receiving substitution of that completed
calculation, not a new proof by differentiating at zero.

At \(d=d_k\), insert the value into the exact (GMF41).
The complete result is
\begin{equation}\tag{GMF45}
\begin{split}
 \boxed{
 \mathcal R_{\rm ar}
 \log\frac{\det Q_{\Fcal,N}[\widetilde m_k]}
              {\det Q_{\Fcal,N}[m_k]}
 =-\frac{\log2}{\pi}I_h k^2+\epsilon_{\Fcal,k}},\\
 \epsilon_{\Fcal,k}=-E_k^{\rm cen}+\mathcal E_{\Fcal,k},
 \qquad
 |\epsilon_{\Fcal,k}|
 \le\mathfrak C_{h,k}^{\rm cen}+4\Omega_{k,d_k}.
\end{split}
\end{equation}
All localization guards (GMF14),(GMF21), the degree
guard \(d_k\le q-2\), and the source theorem's original
finite-amplitude guards remain in force for this receiver.
The separate majority-eigenvalue theorem retains its
own additional divisor-degree guard (GMF34).
On the simple stratum \(q=(k+1)^2\), the coefficient at
order \(q\) in this attained complementary return is
\(-(\log2/\pi)I_h\), since \(k^2/q\to1\) and the displayed
error is \(o_h(q)\).
For fixed \(m_0>1\) the exact displayed scale remains
\(k^2\), with \(q\asymp k^3\); that scale has not been
relabeled as an order-\(q\) coefficient.

% BEGIN COMPLETE FCD SOURCE
% Independent reconstruction of the fixed-d calculation, incoming (36)--(39).
% Include in a document loading amsmath, amssymb, and hyperref.
% This fragment has no external input, graphics, or bibliography dependency.
\section{Exact fixed-degree cumulant determinants with the original source masses}
\label{sec:FCD}

The incoming calculation is the second part of the supplied text headed
\emph{A growing marked flag has large individual metric changes---but an
exponentially small signed return}, equations (36)--(39).
This section proves its finite moment calculation, all four compressed
Hermite traces, the original Gamma constants, and the finite cubic
remainder. The attained-source comparison is retained as an exact receiving
identity at the end. Its analytic error is evaluated by the accompanying
growing-flag localization argument.

\subsection{Original objects and exact coordinate maps}

Retain the actual fixed packet, including its multiplicity:
\[
 h(s)=\prod_{\zeta\in\{\rho,\bar\rho,1-\rho,1-\bar\rho\}}
       (s-\zeta)^{m_0},\qquad
 \rho=\tfrac12+\delta+i\gamma,\quad
 0<\delta<\tfrac12,\quad\gamma>2,\quad h\mid2\xi.
\]
Put \(v_h=2\xi/h\), with its unchanged unit, and
\[
 w_h(y)=\frac{|v_h(1/2+iy)|^2}{2\pi},\qquad
 m_k=w_h^{*k},\qquad c=k/2,\qquad
 q=[1+k(m_0-1)](k+1)^2,\qquad k\equiv1\pmod4.
\]
The original packet gives an even nonnegative source with all finite moments.
Evenness follows from the real coefficients of \(h\) and the conjugation
identity of \(\xi\). Since \(v_h\) is a nonzero analytic function, its
zeros on this line are isolated, so \(w_h\) is positive almost everywhere.
In particular the following original constants are positive:
\begin{equation}
 \begin{split}
 M(\theta)&=\int_{\mathbb R}e^{\theta y}w_h(y)\,dy,\qquad
 M_0=M(0),\qquad M_\alpha=M(\alpha),\qquad \alpha=\pi/2,\\
 V&=\frac1{M_0}\int_{\mathbb R}y^2w_h(y)\,dy>0,\qquad
 D_0=\log(M_\alpha/M_0).
 \end{split}
 \tag{FCD1}\label{eq:FCD1}
\end{equation}
The endpoint moment \(M_\alpha\) is the finite endpoint moment of this
packet; it is not replaced by a probability mass. Let
\[
 \mu_j=M_0^{-1}\int_{\mathbb R}y^jw_h(y)\,dy,\qquad
 c_4=\frac{\mu_4-3V^2}{V^2},\qquad
 c_6=\frac{\mu_6-15\mu_4V+30V^3}{V^3}.
\]
More generally \(c_{2r}=\kappa_{2r}(w_h/M_0)/V^r\), with \(c_2=1\);
the cumulants are the coefficients of the formal logarithm of the moment
series. Only finitely many are used for any fixed degree.

For a positive density \(\nu\), let \(H_d[\nu]\) denote the
\((d+1)\)-by-\((d+1)\) Gram matrix of
\((1,S,\ldots,S^d)\) on the literal line \(S=c+iy\):
\[
 (H_d[\nu])_{ij}
   =\int_{\mathbb R}\overline{(c+iy)^i}(c+iy)^j\nu(y)\,dy,
       \qquad 0\le i,j\le d.
\]
Define the auxiliary probability measure
\[
 d\mu^{(k)}(x)
  =M_0^{-k}\sqrt{kV}\,m_k(\sqrt{kV}\,x)\,dx
\]
and its moment matrix \(G_d^{(k)}\) in \(1,x,\ldots,x^d\).
The exact pullback map on polynomials is
\[
 \mathcal U_k:\mathbb C[S]_{\le d}\longrightarrow\mathbb C[x]_{\le d},
 \qquad P\longmapsto P(c+i\sqrt{kV}\,x).
\]
Its matrix \(B\), with columns indexed by powers of \(S\), has entries
\[
 B_{\ell j}=
 \begin{cases}
 \binom j\ell c^{j-\ell}(i\sqrt{kV})^\ell,&\ell\le j,\\
 0,&\ell>j .
 \end{cases}
\]
The inverse substitutes \(x=(S-c)/(i\sqrt{kV})\).
Consequently, with \(D=d(d+1)/2\),
\begin{equation}
 H_d[m_k]=M_0^kB^*G_d^{(k)}B,\qquad
 \det B=(i\sqrt{kV})^D,\qquad
 \det H_d[m_k]=M_0^{k(d+1)}(kV)^D\det G_d^{(k)} .
 \tag{FCD2}\label{eq:FCD2}
\end{equation}
This equality retains the physical centre, all factors \(i\), the entire
original mass, and the coordinate Jacobian. The auxiliary measure supplies
finite coefficient identities; \(\mathcal U_k\) and \eqref{eq:FCD2}
return them to the unchanged source metric.

\subsection{An exact finite polynomial in \(1/k\)}

For each moment order at most \(2d\), independence in the convolution and
the formal logarithm give
\begin{equation}
 \sum_{n=0}^{2d}\frac{\int x^n\,d\mu^{(k)}(x)}{n!}z^n
 \equiv
 \exp\left(\frac{z^2}{2}
       +\sum_{r=2}^{d}\frac{c_{2r}}{(2r)!}\,
                    k^{1-r}z^{2r}\right)\pmod {z^{2d+1}}.
 \tag{FCD3}\label{eq:FCD3}
\end{equation}
Indeed the cumulant of order \(2r\) of the sum is \(k\) times the
one-packet cumulant, and the exact map \(\mathcal U_k\) contributes
\((kV)^{-r}\). Odd cumulants vanish by evenness. These assertions can
also be obtained by taking the formal logarithm of
\((M(z/\sqrt{kV})/M_0)^k\); no infinite-series approximation is needed.

Write \(\varepsilon=1/k\), and define polynomials \(\mathfrak m_n(\varepsilon)\)
by the finite coefficients in \eqref{eq:FCD3}. Explicitly,
\[
 \mathfrak m_{2a}(\varepsilon)
 =(2a)!\!
 \sum_{\substack{n_1,\ldots,n_a\ge0\\
                         \sum_{r=1}^a r n_r=a}}
 \prod_{r=1}^a
  \frac{(c_{2r}\varepsilon^{r-1})^{n_r}}
       {n_r!((2r)!)^{n_r}},
 \qquad \mathfrak m_{2a+1}=0,\qquad \mathfrak m_0=1.
\]
Here \(c_2=1\) and \(\varepsilon^0=1\), including the \(r=1\) factor.
Define
\begin{equation}
 P_d(\varepsilon)=
 \frac{\det[\mathfrak m_{i+j}(\varepsilon)]_{i,j=0}^{d}}
      {\prod_{j=0}^d j!}.
 \tag{FCD4}\label{eq:FCD4}
\end{equation}
This is a finite real polynomial. At \(\varepsilon=0\) these moments are
Gaussian moments, and the calculation below proves \(P_d(0)=1\).
For the actual positive integer \(k\), its moment matrix is positive
definite: the integral of the square of a nonzero polynomial against
\(\mu^{(k)}\) is strictly positive. In particular \(P_d(1/k)>0\).
Combining \eqref{eq:FCD2} and \eqref{eq:FCD4} proves the exact identity
\begin{equation}
 \boxed{\displaystyle
 \det H_d[m_k]=M_0^{k(d+1)}(kV)^{d(d+1)/2}
                    \left(\prod_{j=0}^d j!\right)P_d(1/k).}
 \tag{FCD5}\label{eq:FCD5}
\end{equation}

\subsection{Hermite identities and all compressed traces}

We use the probabilists' Hermite generating function recorded by
T.~H.~Koornwinder, R.~Wong, R.~Koekoek, R.~F.~Swarttouw and
W.~P.~Reinhardt in the NIST Digital Library of Mathematical Functions,
\href{https://dlmf.nist.gov/18.12.E16}{equation 18.12.16}.
The product formula below is its direct coefficient consequence and agrees
with the probabilists' rescaling of their
\href{https://dlmf.nist.gov/18.18.E23}{equation 18.18.23}.
We give the derivation so that the trace calculation has no unstated
orthogonal-polynomial convention.
Set
\[
 d\gamma(x)=(2\pi)^{-1/2}e^{-x^2/2}\,dx,\qquad
 e_n(x)=\operatorname{He}_n(x)/\sqrt{n!},\qquad
 e^{xz-z^2/2}=\sum_{n\ge0}\operatorname{He}_n(x)z^n/n!.
\]
Completing the square in the Gaussian integral gives
\[
 \int e^{xz-z^2/2}e^{xt-t^2/2}\,d\gamma(x)=e^{zt}.
\]
Comparing the \(z^mt^n\) coefficients proves
\(\int \operatorname{He}_m\operatorname{He}_n\,d\gamma
   =n!\,\mathbf1_{\{m=n\}}\).
The defining generating function also shows that
\(\operatorname{He}_n\) is monic. Its triangular change from monomials has
determinant one; hence the Gaussian monomial Gram determinant equals
\(\prod_{n=0}^d n!\), proving \(P_d(0)=1\).
The identity
\[
 e^{xz-z^2/2}e^{xt-t^2/2}
    =e^{zt}e^{x(z+t)-(z+t)^2/2}
\]
proves
\begin{equation}
 \operatorname{He}_j\operatorname{He}_n
   =\sum_{r=0}^{\min(j,n)}
       r!\binom jr\binom nr\operatorname{He}_{j+n-2r}.
 \tag{FCD6}\label{eq:FCD6}
\end{equation}

Let \(T_j\) be multiplication by \(\operatorname{He}_j\), compressed
orthogonally in \(L^2(\gamma)\) to
\(\operatorname{span}\{e_0,\ldots,e_d\}\). Formula \eqref{eq:FCD6}
proves its complete entries:
\begin{equation}
 (T_j)_{mn}=
 \begin{cases}
 \displaystyle
 \frac{j!\sqrt{m!n!}}{a!\,b!\,r!},
 &
 a=(j+n-m)/2,\quad b=(j+m-n)/2,\quad r=(m+n-j)/2
       \ \text{are nonnegative integers},\\[4pt]
 0,&\text{otherwise}.
 \end{cases}
 \tag{FCD7}\label{eq:FCD7}
\end{equation}
This expression also proves symmetry. Write
\(x^{\underline r}=x(x-1)\cdots(x-r+1)\), with
\(x^{\underline0}=1\). Its value is zero for nonnegative integral
\(x<r\). On the diagonal,
\[
 (T_{2r})_{nn}=\binom{2r}{r}n^{\underline r}.
\]
The polynomial identity
\((n+1)^{\underline{r+1}}-n^{\underline{r+1}}
 =(r+1)n^{\underline r}\)
telescopes. Thus, for every integer \(d\ge0\),
\begin{equation}
 \operatorname{Tr}T_{2r}
    =\frac{\binom{2r}{r}}{r+1}(d+1)^{\underline{r+1}},
 \qquad
 \begin{split}
 \operatorname{Tr}T_4&=2(d+1)d(d-1),\\
 \operatorname{Tr}T_6&=5(d+1)d(d-1)(d-2),\\
 \operatorname{Tr}T_8&=14(d+1)d(d-1)(d-2)(d-3).
 \end{split}
 \tag{FCD8}\label{eq:FCD8}
\end{equation}

All nonzero bands of \(T_4\), before removing indices outside \(0,\ldots,d\),
are
\[
 \begin{split}
 (T_4)_{nn}&=6n(n-1),\\
 (T_4)_{n+2,n}=(T_4)_{n,n+2}
       &=4n\sqrt{(n+1)(n+2)},\\
 (T_4)_{n+4,n}=(T_4)_{n,n+4}
       &=\sqrt{(n+1)(n+2)(n+3)(n+4)} .
 \end{split}
\]
Consequently the compression, including its two boundary bands, gives
\[
 \begin{split}
 \operatorname{Tr}(T_4^2)
 ={}&36\sum_{n=0}^d n^2(n-1)^2
     +32\sum_{n=0}^{d-2}n^2(n+1)(n+2)\\
    &+2\sum_{n=0}^{d-4}(n+1)(n+2)(n+3)(n+4).
 \end{split}
\]
An upper limit below zero denotes the empty sum. Put
\(F_r=(d+1)^{\underline r}\). The polynomial identities
\[
 n^2(n-1)^2=n^{\underline4}+4n^{\underline3}+2n^{\underline2},
 \qquad
 n^2(n+1)(n+2)=(n+2)^{\underline4}+(n+2)^{\underline3}
\]
and the same telescoping formula give, respectively, the three sums
\[
 F_5/5+F_4+2F_3/3,\qquad F_5/5+F_4/4,\qquad F_5/5.
\]
The missing lower-index falling factorials in the shifted sums vanish, so
these formulas include every empty-sum case. Hence
\begin{equation}
 \boxed{\displaystyle
 \operatorname{Tr}(T_4^2)=14F_5+44F_4+24F_3
             =2(d-1)d(d+1)(10-13d+7d^2).}
 \tag{FCD9}\label{eq:FCD9}
\end{equation}
In particular \(T_4=0\) for \(d=0,1\). No large-\(d\) qualification
has been used in \eqref{eq:FCD8} or \eqref{eq:FCD9}.

\subsection{The two exact logarithmic coefficients}

The elementary Gaussian identity
\[
 \int e^{zx}\operatorname{He}_j(x)\,d\gamma(x)
       =z^j e^{z^2/2}
\]
follows by integrating the Hermite generating function against \(e^{zx}\)
and comparing its coefficients. Apply it coefficientwise to
\eqref{eq:FCD3}. For every polynomial \(F\) of degree at most \(2d\),
the first three coefficients of its exact moment functional are
\[
 \begin{split}
 L_\varepsilon(F)
 ={}&\int F\,d\gamma
       +\frac{c_4\varepsilon}{24}
                       \int F\operatorname{He}_4\,d\gamma\\
   &+\varepsilon^2\left[
        \frac{c_6}{720}\int F\operatorname{He}_6\,d\gamma
       +\frac{c_4^2}{1152}\int F\operatorname{He}_8\,d\gamma
                         \right]+\varepsilon^3 R_F(\varepsilon),
 \end{split}
\]
where \(R_F\) is a finite polynomial, possibly zero. This is a
finite coefficient identity of moments. It makes no claim about a
density expansion or a distributional remainder.
In the basis \(e_0,\ldots,e_d\), the Gram matrix is exactly
\[
 I+\varepsilon A+\varepsilon^2 B+\varepsilon^3 R(\varepsilon),
 \qquad
 A=\frac{c_4}{24}T_4,\qquad
 B=\frac{c_6}{720}T_6+\frac{c_4^2}{1152}T_8,
\]
with a polynomial matrix \(R\).
The change from monomials to \(e_n\) has diagonal
\((n!)^{-1/2}\); its squared determinant is
\((\prod_{n=0}^d n!)^{-1}\). Thus this Gram determinant is exactly
\(P_d(\varepsilon)\), including all higher coefficients.

For completeness, the first two logarithmic coefficients are obtained
directly from the determinant permutation formula:
\[
 \det(I+\varepsilon A+\varepsilon^2B+O(\varepsilon^3))
 =1+\varepsilon\operatorname{Tr}A+
 \varepsilon^2\left(\operatorname{Tr}B+
    \tfrac12[(\operatorname{Tr}A)^2-\operatorname{Tr}(A^2)]\right)
 +O(\varepsilon^3).
\]
The diagonal choices give the trace terms; selecting two distinct
first-order entries gives
\(\sum_{i<j}(A_{ii}A_{jj}-A_{ij}A_{ji})\), which is the stated half
difference. Taking the scalar logarithm therefore gives
\[
 [\varepsilon]\log P_d=\operatorname{Tr}A,\qquad
 [\varepsilon^2]\log P_d=\operatorname{Tr}B-\tfrac12\operatorname{Tr}(A^2).
\]
Substitution of \eqref{eq:FCD8} and \eqref{eq:FCD9} yields
\begin{equation}
 \begin{split}
 A_d^{\rm cum}&=\frac{c_4d(d^2-1)}{12},\\
 B_d^{\rm cum}
   &=\frac{d(d^2-1)}{288}
                \left[(16-11d)c_4^2+(2d-4)c_6\right],\\
 \log P_d(1/k)&=\frac{A_d^{\rm cum}}k
                         +\frac{B_d^{\rm cum}}{k^2}
                         +r_{d,k}^{\rm cum}.
 \end{split}
 \tag{FCD10}\label{eq:FCD10}
\end{equation}
The coefficient of \(c_4^2\) contains the trace of the
square of the \emph{compressed} multiplication matrix. Replacing it
by the compression of the full square would discard boundary bands and
would change the answer.
The exact relation between these two operators, on the polynomial domain,
is
\[
 \Pi_d M_{\operatorname{He}_4}^2\Pi_d
       -(\Pi_dM_{\operatorname{He}_4}\Pi_d)^2
   =\Pi_dM_{\operatorname{He}_4}(I-\Pi_d)
                             M_{\operatorname{He}_4}\Pi_d
   =L_d^*L_d,\qquad
 L_d=(I-\Pi_d)M_{\operatorname{He}_4}\Pi_d .
\]
Here \(L_d\) maps the degree-at-most-\(d\) Hermite space into
\(\operatorname{span}\{e_{d+1},\ldots,e_{d+4}\}\).
Inserting \(I=\Pi_d+(I-\Pi_d)\) proves the equality;
self-adjointness of real-polynomial multiplication gives the last
factorization. Formula \eqref{eq:FCD7} evaluates every entry of
this finite boundary map as well.

\subsection{An explicit finite \(k^{-3}\) remainder}

Use the whole polynomial from \eqref{eq:FCD4}, including every
cumulant required through degree \(2d\):
\[
 P_d(\varepsilon)=1+\sum_{j\ge1}a_j\varepsilon^j,\qquad
 C_0=\sum_{j\ge1}|a_j|,\qquad C_3=\sum_{j\ge3}|a_j|.
\]
Both sums are finite. If \(k\ge\max(1,2C_0)\), put
\(u=P_d(1/k)-1\). Then
\[
 |u|\le C_0/k\le\tfrac12,\qquad
 \left|u-\frac{a_1}k-\frac{a_2}{k^2}\right|
       \le C_3/k^3,\qquad
 \left|u-\frac{a_1}k\right|\le C_0/k^2 .
\]
The scalar series for \(\log(1+u)\), absolutely convergent on this
closed half-disc, satisfies
\[
 |\log(1+u)-u+u^2/2|
   \le\sum_{j\ge3}|u|^j/j
   \le\frac{|u|^3}{3(1-|u|)}
   \le\frac{2C_0^3}{3k^3}.
\]
Also
\[
 \frac12\left|u^2-\frac{a_1^2}{k^2}\right|
 \le\frac12\left|u-\frac{a_1}k\right|
                  \left|u+\frac{a_1}k\right|
 \le C_0^2/k^3 .
\]
Since \(A_d^{\rm cum}=a_1\) and
\(B_d^{\rm cum}=a_2-a_1^2/2\), this proves
\begin{equation}
 \boxed{\displaystyle
 |r_{d,k}^{\rm cum}|
 \le \frac{C_3+C_0^2+\tfrac23 C_0^3}{k^3},
 \qquad k\ge\max(1,2C_0).}
 \tag{FCD11}\label{eq:FCD11}
\end{equation}
When \(C_0=0\), \(P_d=1\) and the error is identically zero.
Thus the finite guard includes \(d=0,1\).

\subsection{The original Gamma mass and every monic norm}

The comparison density is literally
\[
 \sigma(y)=\frac{|\Gamma(1/4+iy/2)|^2}{2\pi}.
\]
Here is a derivation of its transform and norms retaining the factor
\(2\pi\). Euler's beta integral
(\href{https://dlmf.nist.gov/5.12.E1}{NIST DLMF 5.12.1})
and \(u=e^x\) give, for \(a>0\),
\[
 \Gamma(a+it)\Gamma(a-it)
   =\Gamma(2a)\int_{\mathbb R}
           e^{itx}(2\cosh(x/2))^{-2a}\,dx.
\]
Fourier inversion therefore gives, initially for real \(z\),
\[
 \int_{\mathbb R}|\Gamma(a+it)|^2e^{-itz}\,dt
        =2\pi\Gamma(2a)(2\cosh(z/2))^{-2a}.
\]
Both sides are analytic for \(|\operatorname{Im}z|<\pi\).
For the integral this follows from the Gamma bound
\(|\Gamma(a+it)|=O_a((1+|t|)^{a-1/2}e^{-\pi|t|/2})\),
recorded in
\href{https://dlmf.nist.gov/5.11.E9}{NIST DLMF 5.11.9};
on compact substrips it also justifies differentiation under the integral.
The analytic identity theorem extends the formula to that strip.
Taking \(z=2i\theta\), \(a=1/4\), and \(y=2t\), gives
\begin{equation}
 \int_{\mathbb R}e^{\theta y}\sigma(y)\,dy
       =\sqrt{2\pi}(\cos\theta)^{-1/2},\qquad
 |\theta|<\pi/2,\qquad M_\sigma=\sqrt{2\pi}.
 \tag{FCD12}\label{eq:FCD12}
\end{equation}

Define the monic polynomials \(p_n\) by
\[
 K(z,y)=\sum_{n\ge0}p_n(y)z^n/n!
       =(1+z^2)^{-1/4}e^{y\arctan z}.
\]
This is the \((\lambda,x,\varphi)=(1/4,y/2,\pi/2)\)
Meixner--Pollaczek generating function in
\href{https://dlmf.nist.gov/18.23.E7}{DLMF 18.23.7}.
Its source note credits M.~E.~H.~Ismail, \emph{Classical and Quantum
Orthogonal Polynomials in One Variable} (2009), equation (5.9.3).
The following calculation proves the needed orthogonality in our own
coordinate and density conventions.
For real \(z,t\) close to zero, \eqref{eq:FCD12} and the cosine addition
formula give
\[
 \int_{\mathbb R}K(z,y)K(t,y)\sigma(y)\,dy
    =M_\sigma(1-zt)^{-1/2}.
\]
Indeed
\(\cos(\arctan z+\arctan t)
 =(1-zt)/\sqrt{(1+z^2)(1+t^2)}\), so all other factors cancel
exactly. The Gamma exponential bound permits coefficientwise
differentiation near zero. Hence
\begin{equation}
 \int p_n(y)p_m(y)\sigma(y)\,dy
       =\mathbf1_{\{n=m\}}M_\sigma n!(1/2)_n .
 \tag{FCD13}\label{eq:FCD13}
\end{equation}
The polynomials \(i^np_n((S-c)/i)\) are monic in \(S\), and their
values at \(c+iy\) are \(i^np_n(y)\). Their monic triangular map
has determinant one; their squared phases have modulus one.
It follows that
\begin{equation}
 \boxed{\displaystyle
 \det H_d[\sigma]
       =\prod_{j=0}^d\left[\sqrt{2\pi}\,j!(1/2)_j\right].}
 \tag{FCD14}\label{eq:FCD14}
\end{equation}

\subsection{Finite source-fill error}

Retain \(b=(\log M)'(\alpha)\), \(T=kb\), and the unchanged
piecewise source
\[
 \widetilde m_k(y)=
 \begin{cases}M_\alpha^k\sigma(y),&|y|\le T,\\
 m_k(y),&|y|>T.
 \end{cases}
\]
Use \(r_*=\tanh(\pi/8)\), \(\theta_*=\pi/4\),
\(\omega_*=\alpha-\theta_*=\pi/4\), and
\(\ell_*=-\log r_*\). These are fixed estimate parameters.
By the power series for \(\arctan z\),
\(|\arctan z|\le\operatorname{artanh}r_*=\pi/8\) on \(|z|=r_*\).
Cauchy's coefficient bound gives
\[
 |p_n(y)|\le n!r_*^{-n}(1-r_*^2)^{-1/4}
                                    e^{(\pi/8)|y|}.
\]
The largest of the \(2n+1\) binomial coefficients
\(\binom{2n}{j}\) is \(\binom{2n}{n}\), so
\(n!/(1/2)_n=4^n/\binom{2n}{n}\le2n+1\).
Expanding a polynomial in the orthogonal basis \eqref{eq:FCD13}
and applying Cauchy--Schwarz thus proves
\[
 |P(c+iy)|^2\le A_d e^{\theta_*|y|}\|P\|_\sigma^2,\qquad
 A_d=\frac{(d+1)^2e^{2d\ell_*}}
             {M_\sigma\sqrt{1-r_*^2}},
       \qquad \deg P\le d .
\]
The sum of \(2n+1\) for \(0\le n\le d\) is \((d+1)^2\);
this proves the displayed constant including its highest-degree factor.
Evenness and the exact convolution transform give
\(\int e^{\alpha|y|}m_k(y)\,dy\le2M_\alpha^k\).
Therefore
\[
 \int_{|y|>T}|P(c+iy)|^2m_k(y)\,dy
  \le2A_dM_\alpha^ke^{-\omega_*T}\|P\|_\sigma^2 .
\]
For the original finite Gamma envelope
\(\sigma(y)\le\overline c_\Gamma
e^{-\alpha|y|}(1+|y|)^{-1/2}\), let
\[
 \eta_{k,d}=\frac{2\overline c_\Gamma A_d}
                 {\omega_*\sqrt{1+T}}e^{-\omega_*T},
 \qquad t_{k,d}=2A_de^{-\omega_*T}.
\]
Integrating the Gamma envelope separately on the two exterior half-lines
proves
\(\int_{|y|>T}|P|^2\sigma\le\eta_{k,d}\|P\|_\sigma^2\).
It follows, for \(\eta_{k,d}<1\), that
\[
 M_\alpha^k(1-\eta_{k,d})H_d[\sigma]
  \preceq H_d[\widetilde m_k]
  \preceq M_\alpha^k(1+t_{k,d})H_d[\sigma].
\]
Define the exact source-fill error
\[
 r_{d,k}^{\rm fill}
   =\log\det H_d[\widetilde m_k]
        -(d+1)k\log M_\alpha-\log\det H_d[\sigma].
\]
Taking the positive eigenvalues of the relative forms proves
\begin{equation}
 \begin{split}
 (d+1)\log(1-\eta_{k,d})
       &\le r_{d,k}^{\rm fill}
                 \le(d+1)\log(1+t_{k,d}),\\
 |r_{d,k}^{\rm fill}|
       &\le(d+1)\{-\log(1-\eta_{k,d})+\log(1+t_{k,d})\}.
 \end{split}
 \tag{FCD15}\label{eq:FCD15}
\end{equation}
For each fixed \(d\), this bound decays exponentially in \(k\).
Strict convexity of the even source transform gives \(b>0\):
its logarithmic second derivative is a positive tilted variance, and
its derivative at zero vanishes.

\subsection{Exact receiver for the attained marked Gram}

Define the explicit expression
\begin{equation}
 \begin{split}
 \mathcal A_{k,d}={}&
 (d+1)kD_0-\frac{d(d+1)}2\log(kV)
 +\frac{d+1}2\log(2\pi)+\sum_{j=0}^d\log(1/2)_j\\
 &-\frac{c_4d(d^2-1)}{12k}
 -\frac{d(d^2-1)}{288k^2}
          [(16-11d)c_4^2+(2d-4)c_6].
 \end{split}
 \tag{FCD16}\label{eq:FCD16}
\end{equation}
All factors \(j!\) cancel between \eqref{eq:FCD5} and
\eqref{eq:FCD14}, whereas \(M_0\), \(M_\alpha\), \(V\),
\(\sqrt{2\pi}\), and \((1/2)_j\) remain exactly as displayed.
The source calculation proves the exact finite identity
\begin{equation}
 \log\frac{\det H_d[\widetilde m_k]}{\det H_d[m_k]}
       =\mathcal A_{k,d}+r_{d,k}^{\rm fill}-r_{d,k}^{\rm cum}.
 \tag{FCD17}\label{eq:FCD17}
\end{equation}

For the original minimum-remainder metric \(G_N[\nu]\), let
\(\mathscr F_{\lambda,d}
 =\operatorname{span}\{v_\lambda,1,S,\ldots,S^d\}\), with
\(v_\lambda=b_\lambda\chi(S)/(S-\lambda)\) and the original
complete-primary coefficient \(b_\lambda\).
For each \(q-1\le N\le2q\), define the exact attained transport difference
\[
 r_{N,k,d}^{\rm att}
  =\log\frac{\det G_N[\widetilde m_k]|_{\mathscr F_{\lambda,d}}}
                  {\det G_N[m_k]|_{\mathscr F_{\lambda,d}}}
    -\log\frac{\det H_d[\widetilde m_k]}{\det H_d[m_k]} .
\]
Thus the receiver is the exact equality
\begin{equation}
 \boxed{\displaystyle
 \log\frac{\det G_N[\widetilde m_k]|_{\mathscr F_{\lambda,d}}}
                  {\det G_N[m_k]|_{\mathscr F_{\lambda,d}}}
  =\mathcal A_{k,d}+r_{N,k,d}^{\rm att}
                          +r_{d,k}^{\rm fill}-r_{d,k}^{\rm cum}.}
 \tag{FCD18}\label{eq:FCD18}
\end{equation}
Equation (GMF26) in the accompanying growing-flag proof evaluates
\(|r_{N,k,d}^{\rm att}|\le\Omega_{k,d}\), with uniform decay in (GMF28),
through the actual relation
minimum, terminal minimum, and their complete Schur cross block.
After that proved estimate is inserted, the complete finite error in
\eqref{eq:FCD18} is
\begin{equation}
 \Omega_{k,d}
 +(d+1)\{-\log(1-\eta_{k,d})+\log(1+t_{k,d})\}
 +\frac{C_3+C_0^2+\tfrac23C_0^3}{k^3}.
 \tag{FCD19}\label{eq:FCD19}
\end{equation}
All finite guards of that localization estimate remain in force, in
addition to \(\eta_{k,d}<1\) and \(k\ge\max(1,2C_0)\).
The quantity \(t_{k,d}\) in this section is exactly \(\zeta_{k,d}\)
in (GMF9), and \eqref{eq:FCD15} is the determinant consequence
of (GMF10), with the source constants independently derived above.
This cumulant calculation alone makes no replacement of an attained
metric by its unquotiented source Gram: their exact difference is
the first error term of \eqref{eq:FCD18}.

The original marked derivative frame has squared change-of-column
determinant
\[
 |b_\lambda|^{2(d+1)}((d+1)!)^2
                  |u|^{-(d+1)(d+2)} .
\]
It is the same coefficient map for both source metrics. Hence it cancels
in their ratio, through the exact rule
\(\det(C^*HC)=|\det C|^2\det H\); it is not assigned the value one.
Consequently \eqref{eq:FCD18} receives the original marked derivative
Gram with precisely the same coefficients and finite error.

\subsection{The first two nontrivial fixed flags}

For \(d=1\), the exact moment matrix is \(\operatorname{diag}(1,1)\).
Thus \(P_1=1\), every cumulant error vanishes, and
\begin{equation}
 \mathcal A_{k,1}=2kD_0-\log(kV)+\log\pi .
 \tag{FCD20}\label{eq:FCD20}
\end{equation}
The exact difference between the attained three-column determinant
ratio and \eqref{eq:FCD20} is
\(r_{N,k,1}^{\rm att}+r_{1,k}^{\rm fill}\).
The accompanying localization estimate makes it exponentially small
uniformly over the original cutoff window.

For \(d=2\), the exact moment matrix and polynomial are
\[
 \begin{pmatrix}
 1&0&1\\
 0&1&0\\
 1&0&3+c_4/k
 \end{pmatrix},
 \qquad P_2(1/k)=1+\frac{c_4}{2k}.
\]
Here \(\prod_{j=0}^2(1/2)_j=3/8\), so
\begin{equation}
 \mathcal A_{k,2}
 =3kD_0-3\log(kV)+\tfrac32\log(2\pi)+\log(3/8)
                 -\frac{c_4}{2k}+\frac{c_4^2}{8k^2}.
 \tag{FCD21}\label{eq:FCD21}
\end{equation}
Equations \eqref{eq:FCD18}--\eqref{eq:FCD19} give its finite remainder.
In this case \(C_0=|c_4|/2\), \(C_3=0\); the direct scalar series also
gives the sharper bound
\[
 |r_{2,k}^{\rm cum}|
    \le\frac{|c_4|^3}{24k^3(1-|c_4|/(2k))}
       \quad\text{for }k>|c_4|/2 .
\]
The sign of the quadratic term in \eqref{eq:FCD21} is positive because
the arithmetic determinant contains \(+\log(1+c_4/(2k))\), which is
subtracted in the source ratio.

\paragraph{Exact-algebra evidence.}
The accompanying \texttt{check\_exact\_cumulants.py} uses rational and
symbolic SymPy arithmetic. It verifies the general polynomial trace
identities, evaluates all three logarithmic coefficients directly from
the original Gaussian moment matrices at each \(d=0,\ldots,20\), and
constructs the entire finite cumulant determinant polynomial for
\(d=0,\ldots,5\). It also checks the finite Gamma norm coefficients.
The run records 95 successful assertions. These are new local checks.
The incoming author's reported Wolfram fixtures were unavailable and
were not executed or verified here. The exact finite proofs above
establish the formulas for all fixed nonnegative integers \(d\);
the finite cases are supplementary checks. None of this fixed-degree
calculation substitutes \(d=d_k\) into a nonuniform remainder or
asserts an absolute observation allocation.

% END COMPLETE FCD SOURCE

\begin{thebibliography}{9}
\bibitem{arrival}
Split-Zero web continuation, \emph{A growing marked flag has
large individual metric changes---but an exponentially small
signed return}, received 19 September 2026, second part of
the supplied attachment, equations (1)--(41).
The finite analytic inequalities and maps are reconstructed
in the present proof.  Its reported external symbolic runs
are not counted as locally executed evidence.

\bibitem{cpt}
Split-Zero mathematical continuation,
\emph{Central-path transfer for the original arithmetic
annihilator metrics}, 19 September 2026,
\texttt{CENTRAL\_PATH\_TRANSFER.tex}, CPT1--CPT43.
The source masses, coordinate map, complete annihilator
minimum correspondence and original Gamma norms are the
ones retained here.

\bibitem{providers}
Split-Zero mathematical continuation,
\texttt{RTM\_GTM\_original\_providers.tex},
RTM1--RTM3 and GTM1, GTM7, preserved in the signed arithmetic
value package of 18 September 2026.

\bibitem{centralProvider}
Split-Zero mathematical continuation,
\texttt{CENTRAL\_SOURCE\_WINDOW.md}, ACW1--ACW13 and
ACW25--ACW26, preserved in the same source package.

\bibitem{markedProvider}
Split-Zero marked arithmetic integration,
\emph{Marked derivatives detect an original central metric
direction}, 19 September 2026,
\texttt{02\_MARKED\_CLASS\_CENTRAL\_SENSITIVITY.md},
B1--B13, preserved in
\texttt{Marked\_Arithmetic\_Integration\_20260919(1).zip}.
The exact partial frame recurrence is retained in (GMF38);
the analytic B7--B11 receivers are derived and strengthened
in (GMF42)--(GMF44).
The exact marked-map receiving proof is
\texttt{MARKED\_ARITHMETIC\_RECOVERY.tex},
MAR1--MAR3 and MAR27--MAR32; its MAR27 is the complete
partial-frame congruence used in (GMF39).

\bibitem{signedCentral}
Split-Zero signed arithmetic continuation,
\emph{Signed arithmetic integrated proof},
\texttt{SIGNED\_ARITHMETIC\_INTEGRATED.tex},
19 September 2026, C1--C32, the original analytic
appendices, and receiving calculations SAI1--SAI12.
The accepted finite-amplitude central endpoint is C32
with its phase correction and finite contour/Robin
error in SAI1--SAI11.

\bibitem{dlmfMPgen}
T.\ H.\ Koornwinder, R.\ Wong, R.\ Koekoek,
R.\ F.\ Swarttouw and W.\ P.\ Reinhardt,
\emph{NIST Digital Library of Mathematical Functions},
Chapter 18, equation 18.23.7.
Its source notes cite M.\ E.\ H.\ Ismail,
\emph{Classical and Quantum Orthogonal Polynomials in One
Variable}, equation (5.9.3).
\url{https://dlmf.nist.gov/18.23.E7}.
Consulted 19 September 2026.

\bibitem{dlmfLegendre}
T.\ H.\ Koornwinder, R.\ Wong, R.\ Koekoek,
R.\ F.\ Swarttouw and W.\ P.\ Reinhardt,
\emph{NIST Digital Library of Mathematical Functions},
Chapter 18, Tables 18.3.1 and 18.9.1,
Legendre norms and three-term recurrence.
\url{https://dlmf.nist.gov/18.3#T1};
\url{https://dlmf.nist.gov/18.9#T1}.
Consulted 19 September 2026.

% The FCD source contains named point-of-use human citations.
\end{thebibliography}
\end{document}
